\scanpage{8}
[Il reste à voir ce qui se passe si $\beta=0$, $\rho=0$ donc $\sigma\sigma'=0$. Posons
\[
\begin{aligned}
a&=\sum_0^n\lambda_i a_i,\quad\text{d'où}\quad \psi(a)=\sum-\lambda_i(\beta_0-\beta_{i-1})a_i,\\
f(a)&=\sum_0^n\lambda_i^2(\beta_0-\beta_{i-1})\unsure{-|+}\bigl(\lambda_0\lambda_1\beta_0+\lambda_1\lambda_2(\beta_0\beta_1)\deleted{+}\ldots+\lambda_{n-1}\lambda_n(\beta_0-\beta_{n-1})\bigr]
\end{aligned}
\]
\mbox{---}\qquad ]

Regardons enfin le noyau de
\[
\operatorname{Aut}(E,f,\pi,\pi')\longrightarrow\operatorname{Aut}(\overbrace{P(V)}^{C})
\]
donc les $u$ du $1^{\mathrm{er}}$ mb qui induisent une homothétie sur $V$. Cette homothétie, $\varepsilon$, conservant $f_V$, doit être dans $\mu_2$, donc $\varepsilon u$ \deleted{conserve} est l'identité sur $V$, donc de la forme
\[
\varepsilon u=\mathrm{id}_E+\deleted{\unsure{\lambda|\varepsilon}}\pi'\otimes\varepsilon a\qquad(a\in E)
\]
\deleted{et \hspace{1em}$\varepsilon\,\mathrm{id}$,}

$\deleted{u=}u=\varepsilon\,\mathrm{id}+\pi'\otimes\deleted{\varepsilon}a$ conserve $f,\pi,\pi'$.

Conserver $\pi,\pi'$
\[
\begin{aligned}
\pi&=\varepsilon\pi+\rho\deleted{\varepsilon}a\qquad\text{i.e.}\quad\boxed{(1-\varepsilon)\pi=\rho\deleted{\varepsilon}a}\\
\pi'&=\varepsilon\pi'+\deleted{\varepsilon}\langle a,\pi'\rangle\pi'\qquad\text{i.e.}\quad\deleted{\boxed{\varepsilon(1+\langle a,\pi'\rangle)=1}}\\
&&\deleted{\unsure{+a|\varepsilon a}}\boxed{\langle a,\pi'\rangle=1-\varepsilon}
\end{aligned}
\]
Enfin, conservation de $f$ \inserted{pour $u$ (ou $\varepsilon u$)} signifie
\[
\boxed{f(a)\pi'+\varepsilon\psi(a)=0}
\]

% archive sheet number: 7

